\section{Apparatus and measurements}\label{sec:apparatus}
The instrument was a scintillator tracker with \InputsDetectorPlanes{} planes and \InputsBarsPerPlane{} live bars per plane. The active width and separation between consecutive layers were determined from the open-sky hit distributions and angular acceptance edge. The geometry used in every forward projection is listed in Table~\ref{tab:apparatus}. Detector dimensions enter in centimetres at the configuration boundary; reconstructed positions and lengths are expressed in metres.

\begin{table}[htbp]
\centering
\caption{Detector geometry and on-site comparison measurements. The ceiling measurements are comparison inputs and do not fix the fitted beam geometry.}\label{tab:apparatus}
\begin{tabular}{lr}
\toprule
Quantity & Value \\
\midrule
Active tracker width & $\InputsActiveWidthCm\,\mathrm{cm}$ \\
Layer separation & $\InputsLayerGapCm\,\mathrm{cm}$ \\
On-site beam-bottom height & $\InputsMeasuredBeamBottom\,\mathrm{m}$ \\
On-site vertical beam depth & $\InputsMeasuredBeamDepth\,\mathrm{m}$ \\
\bottomrule
\end{tabular}
\end{table}

The tracker occupied a reference floor position, pos0, and a displaced position, pos1. An open-sky roof exposure supplied the response reference. The retained exposures contained \DataFirstTracks{} tracks over $\DataFirstHours\,\mathrm{h}$ at pos0, \DataSecondTracks{} tracks over $\DataSecondHours\,\mathrm{h}$ at pos1, and \DataSkyTracks{} tracks over $\DataSkyHours\,\mathrm{h}$ in the sky reference. These counts refer to the angular crop retained by ingestion, rather than all events in the original histograms. Live times were obtained by summing the inter-track interval histogram weighted by the interval-bin centres. Its time unit was inferred from the acquisition rate because the histogram axis was unlabelled.

The world frame places pos0 at the origin, with $z$ upwards and the principal beams running along $y$. The pos1 floor pose was self-calibrated from the transmission data. A rough straight-line separation estimate was not imposed as an exact baseline. The measured beam bottom and depth were retained as independent comparisons; no uncertainty was supplied for those on-site values, so they do not define a formal measurement-consistency test. Fitted heights refer to the configured pos0 reference plane. The direct on-site comparison assumes a common height origin; a separate detector-reference-to-floor offset was not calibrated in this analysis.

\begin{figure}[htbp]
\centering
\includegraphics[width=\linewidth]{setup.pdf}
\caption{Top and side views of the measurement geometry. Detector locations are from the self-calibrated pose; beam boxes show the fitted geometry, not the on-site reference dimensions.}\label{fig:setup}
\end{figure}
